
Donc :
@@M0@@
Or :
@@M1@@
Si on pose @@M2@@, on obtient donc :
@@M3@@
Par ailleurs, @@M4@@; or @@M5@@; donc @@M6@@.

\textbf{Proposition 5.11.} \uline{Avec les notations et les hypothèses de} 5.1, \uline{soit} @@M7@@. \uline{Pour que} @@M8@@, \uline{il faut et il suffit qu'il existe} @@M9@@ \uline{tel que} :
@@M10@@
La démonstration est identique à celle de 5.10, en se plaçant dans les gradués, et en utilisant 5.6.


\section*{6. Etude de la filtration de \texorpdfstring{@@M11@@}{K^{\bullet}(X)}, @@M12@@ ayant un faisceau ample}

Dans ce paragraphe, @@M13@@ est un schéma ayant un faisceau ample @@M14@@. On rappelle que dans ce cas les @@M15@@ définis en termes de faisceaux localement libres de type fini, ou de complexes parfaits, sont égaux (IV, 2.9). On note
